\chapter{Fourier Transform and Momentum-Space Blocks}
\label{ch:fourier-transform-momentum-blocks}

\section{Purpose of this chapter}

After the Jordan--Wigner transformation, the spin chains of interest become
fermionic lattice Hamiltonians.  For open chains, the resulting quadratic
Hamiltonians can be diagonalized directly in real space.  For periodic and
translation-invariant chains, however, the natural variables are momentum modes.
The purpose of this chapter is to fix the Fourier-transform convention used in
these notes and to explain why quadratic Hamiltonians decompose into independent
momentum-space blocks.

There are three convention-dependent points that must be fixed carefully:
\begin{enumerate}[label=(\roman*)]
    \item the sign convention in the Fourier transform;
    \item the allowed momenta, which depend on the fermionic boundary condition;
    \item the harmless but useful global phase that may be included in the
    Fourier transform.
\end{enumerate}
The third point is especially convenient in superconducting chains, where the
pairing gap may be complex.  A uniform phase of the pairing gap can often be moved
between the Hamiltonian and the definition of the momentum-space fermions by a
constant gauge rotation.

\section{Discrete Fourier transform with a boundary twist}

Consider a chain of length $L$ with fermionic operators $c_j,c_j^\dagger$,
$j=1,\ldots,L$, satisfying
\begin{equation}
    \{c_i,c_j^\dagger\}=\delta_{ij},
    \qquad
    \{c_i,c_j\}=\{c_i^\dagger,c_j^\dagger\}=0.
\end{equation}
We impose the general boundary condition
\begin{equation}
\label{eq:general-fermion-boundary-condition}
    c_{j+L}=\ee^{\ii\vartheta}c_j,
    \qquad
    \vartheta\in[0,2\pi).
\end{equation}
The phase $\vartheta$ is a boundary twist.  Periodic boundary conditions
correspond to $\vartheta=0$, while anti-periodic boundary conditions correspond
to $\vartheta=\pi$.

The allowed momenta are the solutions of
\begin{equation}
    \ee^{\ii k L}=\ee^{\ii\vartheta}.
\end{equation}
Thus
\begin{equation}
\label{eq:allowed-momenta-with-twist}
    k_m
    =
    \frac{2\pi m+\vartheta}{L},
    \qquad
    m=0,1,\ldots,L-1,
\end{equation}
where momenta are understood modulo $2\pi$.  We denote this finite set by
\begin{equation}
    \mathcal{K}_\vartheta
    =
    \left\{
    \frac{2\pi m+\vartheta}{L}\;\middle|\;m=0,1,\ldots,L-1
    \right\}.
\end{equation}

\begin{definition}[Fourier transform used in these notes]
\label{def:fourier-transform-with-phase}
For $k\in\mathcal{K}_\vartheta$, we define
\begin{equation}
\label{eq:fourier-transform-with-phase}
    c_k
    =
    \frac{\ee^{-\ii\phi}}{\sqrt{L}}
    \sum_{j=1}^{L}\ee^{-\ii k j}c_j,
    \qquad
    c_j
    =
    \frac{\ee^{\ii\phi}}{\sqrt{L}}
    \sum_{k\in\mathcal{K}_\vartheta}\ee^{\ii k j}c_k.
\end{equation}
Here $\phi$ is a constant Fourier-gauge phase.  Unless otherwise stated, the
default convention is $\phi=0$.
\end{definition}

The phase $\phi$ in \cref{eq:fourier-transform-with-phase} is the phase shown in
the convention
\begin{equation}
    c_k
    =
    \frac{\ee^{-\ii\phi}}{\sqrt{L}}
    \sum_{j=1}^{L}\ee^{-\ii k j}c_j,
    \qquad
    c_j
    =
    \frac{\ee^{\ii\phi}}{\sqrt{L}}
    \sum_{k}\ee^{\ii k j}c_k.
\end{equation}
It is not a boundary twist.  The boundary twist is $\vartheta$, which changes the
allowed values of $k$.  The constant phase $\phi$ only redefines every momentum
operator by the same $U(1)$ phase.

The inverse transform follows from the finite orthogonality relations
\begin{equation}
\label{eq:finite-fourier-orthogonality}
    \frac{1}{L}\sum_{j=1}^{L}\ee^{\ii(k-k')j}
    =
    \delta_{k,k'},
    \qquad
    \frac{1}{L}\sum_{k\in\mathcal{K}_\vartheta}\ee^{\ii k(j-j')}
    =
    \delta_{j,j'} .
\end{equation}
The second identity is valid for $j,j'=1,\ldots,L$.  The common phase $\phi$ drops
out of the orthogonality relations.

\begin{proposition}[Canonical algebra in momentum space]
The operators defined in \cref{eq:fourier-transform-with-phase} obey
\begin{equation}
    \{c_k,c_{k'}^\dagger\}=\delta_{k,k'},
    \qquad
    \{c_k,c_{k'}\}=\{c_k^\dagger,c_{k'}^\dagger\}=0 .
\end{equation}
\end{proposition}

\begin{proof}
Using the real-space canonical anticommutation relations,
\begin{align}
    \{c_k,c_{k'}^\dagger\}
    &=
    \frac{1}{L}
    \sum_{j,j'=1}^{L}
    \ee^{-\ii k j}\ee^{\ii k'j'}
    \{c_j,c_{j'}^\dagger\}
    \nonumber\\
    &=
    \frac{1}{L}\sum_{j=1}^{L}\ee^{\ii(k'-k)j}
    =
    \delta_{k,k'}.
\end{align}
The phases $\ee^{-\ii\phi}$ and $\ee^{\ii\phi}$ cancel.  The other
anticommutators vanish in the same way.
\end{proof}

\begin{remark}[Choice of lattice origin]
The transform above uses $\ee^{\ii k j}$ with $j=1,\ldots,L$.  Some books use
$\ee^{\ii k(j-1)}$ instead.  This is only a $k$-dependent redefinition
$c_k\mapsto \ee^{\ii k}c_k$.  It can shift phase factors between the Fourier
transform and the momentum-space Hamiltonian, but it cannot change spectra,
correlation functions, or topological invariants.  In these notes we keep the
$j=1,\ldots,L$ convention to match the Jordan--Wigner chapters.
\end{remark}

\section{Allowed momenta in Jordan--Wigner parity sectors}

For a periodic spin chain, the Jordan--Wigner transformation does not produce one
single fermionic boundary condition.  Instead, the effective boundary condition
depends on the fermion-parity sector.  With the convention used in these notes,
\cref{eq:effective-jw-boundary-condition} gives
\begin{equation}
    c_{L+1}=-p\,c_1,
    \qquad
    p=(-1)^N=\pm1.
\end{equation}
Therefore
\begin{equation}
\label{eq:jw-parity-momentum-sectors}
\begin{array}{c|c|c|c}
    \text{fermion parity} & p & \text{boundary condition} & \text{allowed momenta} \\
    \hline
    \text{even} & +1 & c_{L+1}=-c_1
    & k_m=\dfrac{(2m+1)\pi}{L} \\
    \text{odd} & -1 & c_{L+1}=c_1
    & k_m=\dfrac{2\pi m}{L}
\end{array}
\end{equation}
with $m=0,1,\ldots,L-1$.

This distinction is negligible for many bulk thermodynamic quantities in the
limit $L\to\infty$, but it is essential for finite-size spectra, exact diagonal
forms, and parity-sensitive ground-state questions.  In particular, periodic
fermions may contain self-conjugate momenta satisfying
\begin{equation}
    k\equiv -k\pmod{2\pi}.
\end{equation}
For even $L$ and periodic boundary conditions, these are $k=0$ and $k=\pi$.  Such
modes must be treated separately in a BdG decomposition because they are not paired
with a distinct momentum $-k$.

In the thermodynamic limit, for a smooth function $f(k)$,
\begin{equation}
\label{eq:thermodynamic-limit-momentum-sum}
    \frac{1}{L}\sum_{k\in\mathcal{K}_\vartheta}f(k)
    \longrightarrow
    \int_{-\pi}^{\pi}\frac{\dd k}{2\pi}\,f(k).
\end{equation}
The boundary angle $\vartheta$ only shifts the mesh by an amount of order $1/L$.

\section{Translation-invariant hopping terms}

The main reason for using Fourier modes is that translation-invariant hopping
terms are diagonal in momentum.  Consider
\begin{equation}
\label{eq:translation-invariant-hopping-real-space}
    H_t
    =
    \sum_{j=1}^{L}\sum_{r}t_r\,c_j^\dagger c_{j+r},
\end{equation}
where the boundary condition is understood and the hopping amplitudes satisfy
\begin{equation}
    t_{-r}=t_r^*
\end{equation}
so that $H_t$ is Hermitian.  Substituting the Fourier transform gives
\begin{align}
    \sum_{j=1}^{L}c_j^\dagger c_{j+r}
    &=
    \sum_{k\in\mathcal{K}_\vartheta}
    \ee^{\ii k r}c_k^\dagger c_k.
\end{align}
Therefore
\begin{equation}
\label{eq:hopping-dispersion-general}
    H_t
    =
    \sum_{k\in\mathcal{K}_\vartheta}\xi(k)c_k^\dagger c_k,
    \qquad
    \xi(k)=\sum_{r}t_r\ee^{\ii k r}.
\end{equation}
The Fourier-gauge phase $\phi$ cancels from all number-conserving bilinears.

\begin{example}[Nearest-neighbour tight-binding chain]
For
\begin{equation}
    H_t=-t\sum_j\left(c_j^\dagger c_{j+1}+c_{j+1}^\dagger c_j\right)
    -\mu\sum_j c_j^\dagger c_j,
\end{equation}
we obtain
\begin{equation}
    H_t=
    \sum_k\xi(k)c_k^\dagger c_k,
    \qquad
    \xi(k)=-\mu-2t\cos k.
\end{equation}
This is the normal-state dispersion that appears in the Kitaev chain and in the
fermionic image of the transverse-field XY chain.
\end{example}

\section{Translation-invariant pairing terms}

Superconducting or Jordan--Wigner-transformed anisotropic spin chains contain
terms that do not conserve fermion number but do conserve fermion parity.  A
convenient translation-invariant pairing Hamiltonian is
\begin{equation}
\label{eq:general-pairing-creation-form}
    H_\Delta
    =
    \frac{1}{2}\sum_{j=1}^{L}\sum_r
    \left(
    \Delta_r c_j^\dagger c_{j+r}^\dagger
    +
    \Delta_r^* c_{j+r}c_j
    \right).
\end{equation}
For spinless fermions, only the antisymmetric part of the pairing is physical,
so one may take
\begin{equation}
    \Delta_{-r}=-\Delta_r.
\end{equation}
The Fourier transform gives
\begin{equation}
\label{eq:momentum-space-pairing-gap-general}
    H_\Delta
    =
    \frac{1}{2}\sum_{k\in\mathcal{K}_\vartheta}
    \left[
    \Delta_\phi(k)c_k^\dagger c_{-k}^\dagger
    +
    \Delta_\phi(k)^*c_{-k}c_k
    \right],
\end{equation}
where
\begin{equation}
\label{eq:gap-function-with-fourier-phase}
    \Delta_\phi(k)
    =
    \ee^{-2\ii\phi}
    \sum_r\Delta_r\ee^{\ii k r}.
\end{equation}
Thus the additional Fourier phase does not affect hopping terms, but it rotates
the superconducting gap by twice the phase.  This factor of two is simply the
charge-two nature of a Cooper-pair operator: a pair creation operator contains two
fermion creation operators.

If the pairing is written instead in the annihilation form
\begin{equation}
    \sum_{j,r}\left(\widetilde{\Delta}_r c_jc_{j+r}
    +\widetilde{\Delta}_r^*c_{j+r}^\dagger c_j^\dagger\right),
\end{equation}
then the coefficient of the annihilation term transforms with $\ee^{+2\ii\phi}$.
The two statements are the same convention expressed from opposite sides of the
Hermitian conjugate pair.

\section{BdG blocks for spinless superconducting chains}

Combining the hopping and pairing terms, define the Nambu spinor
\begin{equation}
\label{eq:nambu-spinor-fourier-chapter}
    \Psi_k
    =
    \begin{pmatrix}
        c_k\\
        c_{-k}^\dagger
    \end{pmatrix}.
\end{equation}
Then a general spinless quadratic Hamiltonian can be written as
\begin{equation}
\label{eq:general-bdg-momentum-form-fourier-chapter}
    H
    =
    \frac{1}{2}\sum_{k\in\mathcal{K}_\vartheta}
    \Psi_k^\dagger\calH_{\rm BdG}(k)\Psi_k
    +E_0,
\end{equation}
with
\begin{equation}
\label{eq:bdg-block-general-fourier-chapter}
    \calH_{\rm BdG}(k)
    =
    \begin{pmatrix}
        \xi(k) & \Delta_\phi(k)\\
        \Delta_\phi(k)^* & -\xi(-k)
    \end{pmatrix}.
\end{equation}
For most nearest-neighbour chains considered in these notes, $\xi(-k)=\xi(k)$ and
$\Delta_\phi(-k)=-\Delta_\phi(k)$.

Equivalently, after dropping any term proportional to the identity matrix, the BdG block can be
written as
\begin{equation}
\label{eq:bdg-d-vector-fourier-chapter}
    \calH_{\rm BdG}(k)
    =
    d_x(k)\tau_x+d_y(k)\tau_y+d_z(k)\tau_z,
\end{equation}
with
\begin{equation}
    d_x(k)=\operatorname{Re}\Delta_\phi(k),
    \qquad
    d_y(k)=-\operatorname{Im}\Delta_\phi(k),
    \qquad
    d_z(k)=\xi(k).
\end{equation}
The quasiparticle energies are
\begin{equation}
\label{eq:bdg-spectrum-fourier-chapter}
    E_\pm(k)=\pm\sqrt{\xi(k)^2+\abs{\Delta_\phi(k)}^2}.
\end{equation}
Hence a constant Fourier-gauge phase changes the orientation of the vector
$(d_x,d_y)$ in the pairing plane but does not change the spectrum.

\section{Independent momentum-space blocks}

A translation-invariant number-conserving quadratic Hamiltonian has the form
\begin{equation}
    H=\sum_k\xi(k)c_k^\dagger c_k.
\end{equation}
It is already decomposed into independent one-mode blocks labelled by $k$.

A pairing Hamiltonian is different: it couples $k$ and $-k$.  Therefore its basic
blocks are not single momenta but pairs $(k,-k)$.  Let $\mathcal{K}_+$ be a choice
of one representative from each pair $\{k,-k\}$, excluding self-conjugate momenta.
Then
\begin{equation}
\label{eq:paired-momentum-block-decomposition}
    H
    =
    \sum_{k\in\mathcal{K}_+}H_{k,-k}
    +H_{\rm self}
    +E_0.
\end{equation}
Here $H_{\rm self}$ contains possible self-conjugate modes such as $k=0$ and
$k=\pi$ under periodic boundary conditions.  For odd-parity spinless pairing,
$\Delta_\phi(k)=0$ at self-conjugate momenta, so these modes remain unpaired.

For a non-self-conjugate pair, the even-parity local Hilbert space is spanned by
\begin{equation}
    \ket{0}_{k,-k},
    \qquad
    c_k^\dagger c_{-k}^\dagger\ket{0}_{k,-k}.
\end{equation}
In this subspace, the BdG block describes a two-level system.  This is why the
transverse-field Ising chain, the transverse-field XY chain, and the Kitaev chain
all reduce to products of independent two-level problems in momentum space.

\begin{principle}[Momentum-space block principle]
A translation-invariant quadratic fermion Hamiltonian decomposes into independent
momentum blocks.  Number-conserving Hamiltonians decompose into $k$ blocks.
Spinless superconducting Hamiltonians decompose into $(k,-k)$ BdG blocks, plus
possible self-conjugate modes.
\end{principle}

\section{Physical meaning of the extra Fourier phase}

The constant phase $\phi$ in \cref{eq:fourier-transform-with-phase} implements the
operator redefinition
\begin{equation}
    c_k^{(\phi)}=\ee^{-\ii\phi}c_k^{(0)}.
\end{equation}
In Nambu space,
\begin{equation}
    \Psi_k^{(\phi)}
    =
    \begin{pmatrix}
        \ee^{-\ii\phi} & 0\\
        0 & \ee^{\ii\phi}
    \end{pmatrix}
    \Psi_k^{(0)}
    =
    \ee^{-\ii\phi\tau_z}\Psi_k^{(0)}.
\end{equation}
Therefore BdG matrices in the two conventions are related by the unitary gauge
transformation
\begin{equation}
\label{eq:bdg-gauge-transformation-fourier-phase}
    \calH_{\rm BdG}^{(\phi)}(k)
    =
    \ee^{-\ii\phi\tau_z}
    \calH_{\rm BdG}^{(0)}(k)
    \ee^{\ii\phi\tau_z}.
\end{equation}
This rotates the $\tau_x$--$\tau_y$ components of the BdG vector by a constant
angle and leaves $d_z(k)$ invariant.

The distinction between $\vartheta$ and $\phi$ is important:
\begin{itemize}
    \item $\vartheta$ changes the boundary condition and therefore changes the
    allowed momenta at finite $L$;
    \item $\phi$ is a uniform gauge choice for momentum-space fermions and does
    not change the allowed momenta.
\end{itemize}
A boundary twist can represent a physical flux through a ring.  A constant
Fourier-gauge phase by itself cannot.

\section{Example: complex pairing phase in the Kitaev chain}

Consider the Kitaev chain written in the creation-pair convention
\begin{equation}
\label{eq:kitaev-chain-complex-gap-fourier-chapter}
    H_{\rm K}
    =
    -\mu\sum_j\left(c_j^\dagger c_j-\frac{1}{2}\right)
    -t\sum_j\left(c_j^\dagger c_{j+1}+c_{j+1}^\dagger c_j\right)
    +\sum_j\left(
        \Delta c_j^\dagger c_{j+1}^\dagger
        +
        \Delta^* c_{j+1}c_j
    \right),
\end{equation}
where
\begin{equation}
    \Delta=\abs{\Delta}\ee^{\ii\theta_\Delta}.
\end{equation}
Using \cref{eq:gap-function-with-fourier-phase}, the off-diagonal BdG gap is
\begin{equation}
\label{eq:kitaev-gap-with-fourier-phase}
    \Delta_\phi(k)
    =
    2\ii\abs{\Delta}\,
    \ee^{\ii(\theta_\Delta-2\phi)}\sin k.
\end{equation}
The normal dispersion is
\begin{equation}
    \xi(k)=-\mu-2t\cos k.
\end{equation}
Thus
\begin{equation}
\label{eq:kitaev-bdg-complex-gap-fourier-chapter}
    \calH_{\rm K}(k)
    =
    \begin{pmatrix}
        \xi(k) & 2\ii\abs{\Delta}\ee^{\ii(\theta_\Delta-2\phi)}\sin k\\
        -2\ii\abs{\Delta}\ee^{-\ii(\theta_\Delta-2\phi)}\sin k & -\xi(k)
    \end{pmatrix}.
\end{equation}
The spectrum is
\begin{equation}
    E_\pm(k)
    =
    \pm\sqrt{\left(\mu+2t\cos k\right)^2
    +4\abs{\Delta}^2\sin^2k},
\end{equation}
which is independent of $\theta_\Delta$ and $\phi$.

Equation \cref{eq:kitaev-gap-with-fourier-phase} shows the role of the extra
Fourier phase.  Choosing
\begin{equation}
    \phi=\frac{\theta_\Delta}{2}
\end{equation}
removes the explicit superconducting phase from the gap and gives the standard
purely imaginary $p$-wave form
\begin{equation}
    \Delta_\phi(k)=2\ii\abs{\Delta}\sin k.
\end{equation}
Alternatively, one may choose a different $\phi$ to make the off-diagonal gap point
along another fixed direction in the $\tau_x$--$\tau_y$ plane.

For an isolated uniform Kitaev chain, the global phase $\theta_\Delta$ is therefore
not an independent bulk observable.  It is a superconducting gauge phase.  The
bulk spectrum and the topological phase boundary
\begin{equation}
    \abs{\mu}=2\abs{t}
\end{equation}
are unchanged by it.  What is physical is a gauge-invariant relative phase: for
example, the phase difference between two superconductors in a Josephson junction,
a spatial gradient of the superconducting phase, or a boundary twist that cannot
be removed globally.  In such situations the phase is tied to current, flux, or
junction physics and cannot be discarded as a mere convention.

From the BdG-vector viewpoint, a uniform complex phase of $\Delta$ rotates the
closed curve
\begin{equation}
    k\mapsto (d_x(k),d_y(k),d_z(k))
\end{equation}
around the $d_z$ axis.  This does not change whether the curve winds around the
origin in the pairing plane.  Hence the topological classification of the uniform
Kitaev chain is unaffected by the constant gap phase, even though the explicit
matrix representation of the BdG block changes.

\section{Summary of conventions}

The Fourier convention used in the rest of these notes is
\begin{equation}
    c_k
    =
    \frac{\ee^{-\ii\phi}}{\sqrt{L}}
    \sum_{j=1}^{L}\ee^{-\ii k j}c_j,
    \qquad
    c_j
    =
    \frac{\ee^{\ii\phi}}{\sqrt{L}}
    \sum_{k\in\mathcal{K}_\vartheta}\ee^{\ii k j}c_k,
\end{equation}
with default $\phi=0$.  The allowed momenta are
\begin{equation}
    k_m=\frac{2\pi m+\vartheta}{L},
    \qquad
    m=0,1,\ldots,L-1.
\end{equation}
For Jordan--Wigner fermions obtained from a periodic spin chain,
\begin{equation}
\begin{array}{c|c|c}
    \text{sector} & \text{boundary condition} & \text{momenta} \\
    \hline
    (-1)^N=+1 & \text{anti-periodic} & k_m=\dfrac{(2m+1)\pi}{L} \\
    (-1)^N=-1 & \text{periodic} & k_m=\dfrac{2\pi m}{L}
\end{array}
\end{equation}
Number-conserving hopping terms give one-mode blocks labelled by $k$.  Pairing
terms give BdG blocks labelled by $(k,-k)$.  The extra phase $\phi$ is a useful
Fourier-gauge convention: it cancels from hopping terms but shifts the phase of a
pairing gap by twice its value.
